\section{Related Work}

\paragraph{Differentiable Optimization and Decision-Focused Learning.}

Differentiable optimization integrates structured optimization problems as layers within neural networks. OptNet~\cite{amos2017optnet} introduced differentiable quadratic programming layers via implicit differentiation of KKT systems. Subsequent work extended these ideas to linear, conic, and general convex programs~\cite{agrawal2019differentiable, gould2016deep}. These approaches compute gradients of the optimal solution map by differentiating through the KKT conditions, typically requiring the solution of linear systems involving both primal and dual variables.

Decision-focused learning (DFL) builds on this paradigm by training predictive models to optimize downstream decision quality rather than prediction accuracy. The ``Smart Predict-Then-Optimize'' framework~\cite{elmachtoub2022smart} formalized this perspective, showing that minimizing prediction error can be misaligned with decision performance, particularly in linear and combinatorial optimization problems. Subsequent work has extended DFL to structured prediction, combinatorial solvers, and control settings, often formulating the problem as a bilevel optimization task where the upper level learns predictive parameters and the lower level solves an optimization problem.

Most existing DFL approaches rely on general-purpose KKT-based differentiation. While broadly applicable, these methods require solving linear systems whose dimension scales with the total number of primal and dual variables, which can introduce significant computational overhead and numerical sensitivity in structured control problems.

In contrast, we exploit the occupancy-measure linear programming formulation of finite discounted MDPs. Leveraging the non-degeneracy of discounted occupancy polytopes under strictly positive initial distributions, we derive an exact local gradient expression that depends only on the active basis matrix of dimension $|\mathcal{S}| \times |\mathcal{S}|$. This yields a structurally tailored and computationally efficient DFL gradient specialized to finite discounted MDPs, rather than relying on general KKT differentiation.

\paragraph{State Aggregation in Markov Decision Processes.}

State aggregation is a classical technique in approximate dynamic programming and reinforcement learning for reducing the dimensionality of large MDPs~\cite{bertsekas1996neuro, li2006towards}. Aggregation constructs an abstract MDP by partitioning the state space, trading representational fidelity for computational tractability. Prior work has characterized error bounds induced by aggregation and studied conditions under which optimal or near-optimal policies are preserved~\cite{singh1995reinforcement, li2006towards}.

In linear programming formulations of MDPs, aggregation can be interpreted as projecting the occupancy polytope onto a reduced subspace defined by the aggregation mapping. This reduces both the number of constraints and the dimensionality of the feasible region. Consequently, it also shrinks the size of the active basis matrix involved in differentiation.

Our work connects classical state aggregation with differentiable occupancy-based DFL. Beyond reducing computational complexity, aggregation alters the geometry of the occupancy polytope and can mitigate sensitivity to small parameter perturbations that induce basis changes. This provides both scalability and numerical stability benefits when differentiating through LP-based MDP solvers.